\documentclass[a4paper,12pt]{article}

\usepackage[utf8]{inputenc}
\usepackage[T1]{polski}
\usepackage{helvet}
\usepackage{graphicx}
\usepackage{color}
\usepackage{geometry}

\geometry{hmargin={2cm, 2cm}, height=10.0in}

\begin{document}

% =====  STRONA TYTULOWA PRACY INŻYNIERSKIEJ ====
% ostatnia modyfikacja: 2011/03/09, K. Malarz

\thispagestyle{empty}

%% ------------------------ NAGLOWEK STRONY ---------------------------------
\includegraphics[height=37.5mm]{agh_nzw_a_pl_1w_wbr_cmyk}\\
\rule{30mm}{0pt}
{\large\textsf{Wydział Fizyki i Informatyki Stosowanej}}\\
\rule{\textwidth}{3pt}\\
\rule[2ex]
{\textwidth}{1pt}\\
\vspace{7ex}
\begin{center}
{\bf\LARGE\textsf{Praca inżynierska}}\\
\vspace{13ex}
% --------------------------- IMIE I NAZWISKO -------------------------------
{\bf\Large\textsf{Patryk Chodur}}\\
\vspace{3ex}
{\sf \small kierunek studiów:} {\bf\small\textsf{informatyka stosowana}}\\
\vspace{15ex}
%% ------------------------ TYTUL PRACY --------------------------------------
{\bf\huge\textsf{Opracowanie wtyczki do dekodowania
	pakietów Ethernet z~wykorzystaniem oprogramowania Wireshark
}}\\
\vspace{14ex}
%% ------------------------ OPIEKUN PRACY ------------------------------------
{\sf \Large Opiekun:} {\bf\Large\textsf{dr hab. inż. Bartosz Mindur}}\\
\vspace{22ex}
\textsf{\bf\large\textsf{Kraków, styczeń 2021}}
\end{center}
%% =====  STRONA TYTUŁOWA PRACY INŻYNIERSKIEJ  ====

\newpage

%% =====  TYŁ STRONY TYTUŁOWEJ PRACY INŻYNIERSKIEJ  ====
\begin{center}
	{\bf\large\textsf{Oświadczenie studenta}}
\end{center}


{\sf Uprzedzony(-a) o odpowiedzialności karnej na podstawie art. 115 ust. 1 i 2 ustawy z dnia 4 lutego 1994 r. o prawie autorskim i prawach pokrewnych (t.j. Dz. U. z 2018 r. poz. 1191 z późn. zm.): ,,Kto przywłaszcza sobie autorstwo albo wprowadza w błąd co do autorstwa całości lub części cudzego utworu albo artystycznego wykonania, podlega grzywnie, karze ograniczenia wolności albo pozbawienia wolności do lat 3. Tej samej karze podlega, kto rozpowszechnia bez podania nazwiska lub pseudonimu twórcy cudzy utwór w wersji oryginalnej albo w postaci opracowania, artystyczne wykonanie albo publicznie zniekształca taki utwór, artystyczne wykonanie, fonogram, wideogram lub nadanie.'', a także uprzedzony(-a) o odpowiedzialności dyscyplinarnej na podstawie art. 307 ust. 1 ustawy z dnia 20 lipca 2018 r. Prawo o szkolnictwie wyższym i nauce (Dz. U. z 2018 r. poz. 1668 z późn. zm.) ,,Student podlega odpowiedzialności dyscyplinarnej za naruszenie przepisów obowiązujących w uczelni oraz za czyn uchybiający godności studenta.'', oświadczam, że niniejszą pracę dyplomową wykonałem(-am) osobiście i samodzielnie i nie korzystałem(-am) ze źródeł innych niż wymienione w pracy.

\bigskip

Jednocześnie Uczelnia informuje, że zgodnie z art. 15a ww. ustawy o prawie autorskim i prawach pokrewnych Uczelni przysługuje pierwszeństwo w opublikowaniu pracy dyplomowej studenta. Jeżeli Uczelnia nie opublikowała pracy dyplomowej w terminie 6 miesięcy od dnia jej obrony, autor może ją opublikować, chyba że praca jest częścią utworu zbiorowego. Ponadto Uczelnia jako podmiot, o którym mowa w art. 7 ust. 1 pkt 1 ustawy z dnia 20 lipca 2018 r. --- Prawo o szkolnictwie wyższym i nauce (Dz. U. z 2018 r. poz. 1668 z późn. zm.), może korzystać bez wynagrodzenia i bez konieczności uzyskania zgody autora z utworu stworzonego przez studenta w wyniku wykonywania obowiązków związanych z odbywaniem studiów, udostępniać utwór ministrowi właściwemu do spraw szkolnictwa wyższego i nauki oraz korzystać z utworów znajdujących się w prowadzonych przez niego bazach danych, w celu sprawdzania z wykorzystaniem systemu antyplagiatowego. Minister właściwy do spraw szkolnictwa wyższego i nauki może korzystać z prac dyplomowych znajdujących się w prowadzonych przez niego bazach danych w zakresie niezbędnym do zapewnienia prawidłowego utrzymania i rozwoju tych baz oraz współpracujących z nimi systemów informatycznych.}

\vspace{14ex}

\begin{center}
\begin{tabular}{lr}
~~~~~~~~~~~~~~~~~~~~~~~~~~~~~~~~~~~~~~~~~~~~~~~~~~~~~~~~~~~~~~~~~ &
................................................................. \\
~ & {\sf (czytelny podpis)} \\
\end{tabular}
\end{center}

%% =====  TYL STRONY TYTULOWEJ PRACY INŻYNIERSKIEJ  ====

\newpage
\linespread{1.3}
\selectfont

\noindent
Na kolejnych dwóch stronach proszę dołączyć kolejno recenzje pracy popełnione przez Opiekuna oraz Recenzenta (wydrukowane z systemu MISIO i~podpisane przez odpowiednio Opiekuna i~Recenzenta pracy). Papierową wersję pracy (zawierającą podpisane recenzje) proszę złożyć w~dziekanacie celem rejestracji.

Wstęp

Wireshark jest oprogramowaniem do analizy ruchu internetowego na poziomie pakietów wydanym na otwartej licencji GPL2

Cel pracy

Celem pracy jest wykonanie dekodera pakietów bazujących na protokole UDP i~przenoszących konkretne dane pomiarowe.
Dekoder napisany jest jako wtyczka do programu Wireshark. Podejście takie umożliwia szybszą rekompilację wtyczki,
niewymagającą przekomilowywania całego programu, a~także użycie gotowej wtyczki z~dowolną instalacją programu Wireshark,
o~ile zachowana jest zgodność systemu operacyjnego i~wersji programu Wireshark. Instalacja polega na przekopiowaniu
pojedyńczego pliku do katalogu wtyczek programu. Alternatywą dla takiego rozwiązania jest stworzenie wbudowanego dekodera,
ale ze względu na liczne wady - m.~in.: dłuższy czas rekompilacji oraz konieczność ingerencji w~kod źródłowy programu - rozwiązanie
to nie zostało zastosowane.


Narzędzia

Wtyczka napisana została w języku C, bez jawnego użycia zewnętrznych bibliotek. Wykorzystane pliki nagłówkowe
pochodzą ze standardowej biblioteki C, oraz kodu źródłowego programu Wireshark. Program Wireshark używa wielu
bibliotek, takich jak qt oraz zlib, ale na wyszczególnienie zasługuje biblioteka glib w wersji drugiej, której
elementy zostały użyte w kodzie źródłowym aplikacji. Glib2 dostarcza projektowi standardowych typów o stałym
rozmiarze wykorzystywanych przez niektóre funkcje Wireshark.

Cały kodu projektu napisany został w programie Vim.

Podczas kompilacji projektu wykorzystywane jest narzędzie cmake. Służy ono do zarządzania procesem kompilacji,
umożliwiając jednolity proces budowy projektu
na każdym ze wspieranych systemów operacyjnych. Wykrywa ono konieczne do zbudowania narzędzia, informując
o brakujących, a także, w zależności od systemu operacyjnego, tworzy skrypty kompilacyjne. Na systemach Unixowych
oznacza to zwykle generowanie plików Makefile, na systemie Windows - projektu w Visual Studio. Użycie cmake'a
jest pierwszym krokiem standardowej kompilacji programu Wireshark oraz jego wtyczek.

Do kompilacji wymagany jest kompilator do języka C. Decyzję o wykorzystaniu konkretnego narzędzia podejmuje cmake.
W razie braku stosownego kompilatora cmake poinformuje użytkownika o problemie. Ponadto wymagany jest szereg
bibliotek, które zależą od używanego systemu operacyjnego.

Podczas tworzenia wtyczki wykorzystane zostały dwa proste narzędzia napisane przez autora pracy:
- printhex - program napisany w C służący do wyświetlania otrzymanych bajtów w zapisie szesnastkowym. Standardowo
czyta on z stdin i wysyła output na stdout, lecz możliwe jest podanie nazwy plików wejściowego oraz wyjściowego
- reverse - bardzo prosty program zamieniający kolejność bitów w masce - np 0x0FCull na 0x3F00000000000000ull.
Nie posiada on żadnych dodatkowych funkcjonalności i czyta ze stdin.





Całość kodu pisana była w programie Vim





\vspace{85mm}


\end{document}

